\begin{figure*}[t]
\centering\color{revisedred}
\begin{tikzpicture}[
 foundation/.style={draw=teal!70,fill=wash,rounded corners=2pt,align=center,text width=4.8cm,minimum height=1.25cm,inner sep=5pt,font=\small},
 layer/.style={draw=teal!70,fill=white,rounded corners=2pt,align=center,text width=16.2cm,minimum height=0.85cm,inner sep=5pt,font=\small},
 component/.style={draw=teal!70,fill=wash,rounded corners=2pt,align=center,text width=6.8cm,minimum height=1.1cm,inner sep=5pt,font=\small},
 connection/.style={-{Latex[length=2mm]},draw=teal,thick}]
\node[foundation] (game) at (0,0) {\textit{Game Theory}\\Agents, preference weights, actions\\and interdependent outcomes};
\node[foundation] (graph) at (5.7,0) {\textit{Social Graphs}\\Agent swarm communication, directed\\access and message timing};
\node[foundation] (prob) at (11.4,0) {\textit{Probability and Statistics}\\Declared uncertainty, sampling\\and outcome-based evaluation};
\node[layer] (world) at (5.7,-2) {\textit{Boyun} (inspired by \textit{MiroFish}): reviewed world and \textit{agent swarm} specification\\Swarm of agents, resources, social-graph observations, events and settlement rules};
\draw[connection] (game.south)--(game.south |- world.north);
\draw[connection] (graph.south)--(world.north);
\draw[connection] (prob.south)--(prob.south |- world.north);
\node[component] (laya) at (1.4,-3.9) {Local agent decision model: \textit{Laya}\\Permitted swarm observations $\rightarrow$ option scores and actions};
\node[component] (engine) at (10,-3.9) {Execution ledgers and social graph\\Selected agent actions $\rightarrow$ message delivery, price and state changes};
\draw[connection] (world.south -| laya.north)--(laya.north);
\draw[connection] (world.south -| engine.north)--(engine.north);
\draw[connection] (laya.east)--(engine.west);
\node[layer] (record) at (5.7,-5.7) {Recorded swarm trajectories, price paths and scenario analysis\\Inspect agent choices, social-graph propagation and consequences; replay execution};
\draw[connection] (engine.south)--(engine.south |- record.north);
\end{tikzpicture}
\caption{\rev{Foundations of \textit{Boyun}, inspired by \textit{MiroFish}. The
upper row identifies conceptual foundations; the lower layers show how modeled
actors operate as a \textit{swarm of agents} communicating over a \textit{social
graph} while self-hosted \textit{Laya} supplies live agent decisions locally and
deterministic ledgers apply the specified rules.} Probability and statistical tools are separate
analytical components, not automatic calibration of the agent decision model.
Probability calculations use declared inputs; neither the diagram nor replay
implies calibrated predictions or empirically validated social influence.}
\label{fig:foundations}
\end{figure*}

\section{Introduction}\label{sec:introduction}
Financial and economic decisions depend on how people and organizations
respond to one another under uncertainty. A protocol designer evaluating an
incentive must weigh its direct cost, whether partners join, and how other
actors react, since those reactions alter the choices available later.
Comparing alternatives therefore requires scenarios that make actors,
constraints, and assumed events explicit while showing how their interactions
produce different outcomes.

To that end, \textit{agent-based modeling} represents a system through
\rev{an interacting \textit{swarm of agents}---autonomous decision-makers whose
local observations and actions follow explicit interaction rules---}so its
conclusions depend on both the assumptions and their implementation
\cite{grimm2020,galan2009}. However, although language-model agents can
condition choices on context \cite{park2023,yang2025}, an expressed choice
alone does not specify its economic consequences. For example, a disclosure
decision \rev{emitted by one agent in a swarm} still requires rules defining
which \rev{connected agents receive the message over a social graph}, when they
react, and what resource or state changes follow.

Along another dimension, social-simulation frameworks examine \rev{how an
\textit{agent swarm} interacts across a shared social graph}. \textit{OASIS}
(Open Agent Social Interaction Simulations) provides one such platform, while
\textit{MiroFish} uses \rev{an agent-swarm social-simulation} workflow to
explore scenarios \cite{yang2025,mirofish}. \rev{Inspired directly by
\textit{MiroFish}'s paradigm of constructing worlds, orchestrating an
interacting swarm of agents, and synthesizing scenario reports,} \textit{Boyun}
pursues specialized financial and economic analysis, combining complementary
methods to connect strategic incentives, \rev{social-graph communication across
an agent swarm,} and numerical consequences.

\rev{Building on that \textit{MiroFish}-inspired agent-swarm foundation,} we
present \textit{Boyun} ({\chinesefont 拨云}, {\chinesefont bō yún}), a
self-hosted \rev{\textit{agent-swarm}} framework for exploring decisions
through explicit simulated worlds. Its Chinese name means ``part the clouds''
and draws on {\chinesefont 拨云见日}, in which clouds are parted to reveal the
sun. \rev{In \textit{Boyun}, market and protocol actors are modeled as
autonomous \textit{agents} that communicate with one another across a
\textit{social graph} as a \textit{swarm of agents}.} The name expresses the
aim of bringing \rev{that agent swarm}, its incentives, assumptions,
\rev{social-graph relationships}, and possible consequences into view before
action, while keeping that view conditional on the evidence and assumptions
used to construct it. Figure~\ref{fig:foundations} summarizes the foundations
and their operational roles.

First, to give this \rev{agent swarm} an executable structure, \textit{Boyun}
places \textit{game theory} at the center of world construction. \rev{Agents in
the swarm} have preference weights, resources, information, and feasible
actions, and the world's rules determine how their combined choices affect
outcomes \cite{polak2,nash1950}. This structure makes incentives, coordination,
and conflicts of interest explicit without assuming that model-selected actions
are optimal or form an equilibrium. The distinction matters whenever
\rev{agents} value different outcomes: a founder \rev{agent} may prioritize
participation, whereas a partner \rev{agent} may prioritize credibility or
retained capital. Such differences must alter feasible actions and their
consequences rather than remain descriptive labels in a persona.

On the other hand, incentives alone do not specify who knows what when \rev{inside
the swarm}. \textit{Boyun} therefore \rev{connects the \textit{swarm of agents}
through a \textit{social graph}---formalized as a \textit{directed information
graph}---that specifies which agents receive another agent's messages and how
signals propagate across the swarm}. Message availability and timing thus
become explicit parts of the world rather than hidden assumptions. Even so,
these \rev{social-graph} links describe permitted delivery rather than
empirically measured social influence, allowing an analyst to test how a
disclosure, delayed message, or restricted channel changes \rev{what each agent
in the swarm observes} without giving every \rev{agent} the author's complete
view.

Apart from strategic incentives and \rev{social-graph} communication, events
and outcomes often remain uncertain. \textit{Probability and statistics}
provide a third foundation: \textit{Bayesian} calculations update supplied
beliefs, seeded sampling explores declared alternatives, and proper scoring
rules evaluate probabilistic predictions \cite{pishronik2014,brier1950,gneiting2007}.
Nevertheless, these tools expose uncertainty only under stated assumptions;
neither repeated simulation nor model confidence establishes real-world odds.

Once those three foundations define a world, the engine must obtain \rev{choices
across the \textit{agent swarm}} while keeping local inference and recorded
execution explicit. Because unpublished launch plans, treasury limits, and risk
parameters are sensitive, \textit{Boyun} focuses on private, self-hosted
execution using an open-weight \rev{\textit{agent decision model}} that runs
locally on the user's hardware. Each \rev{agent in the swarm} receives only its
permitted \rev{social-graph and ledger} observation, and \textit{Laya}
\cite{layamodel}---the local \rev{agent} model used in the reported case
studies---returns a structured choice and normalized \textit{option scores}
over the \rev{agent's} \textit{affordable action menu}. Alternative open-weight
models are a future evaluation direction rather than a capability demonstrated
by these runs. Crucially, option scores express model-assigned weights rather
than validated real-world probabilities. Deterministic strategic, market, and
lending ledgers then settle the chosen actions, \rev{route messages across the
social graph,} enforce conservation, and save every observation and choice so
runs replay exactly without repeating inference, while leaving the model's
internal deliberation unobserved.

Moreover, self-hosted inference addresses a practical budget constraint.
\rev{Evaluating an \textit{agent swarm} across many decision points, social-graph
interactions,} and matched policy branches requires tens or hundreds of
thousands of model queries, which quickly exhaust small hosted-API budgets while
sending private scenario states off the user's machine. Running \textit{Laya}
locally eliminates per-call API fees and keeps all \rev{swarm} observations
on-device, though operational efficiency and behavioral accuracy remain separate
evaluation questions.

Given these capabilities and boundaries, the paper asks: \emph{Can a
self-hosted, explicitly specified \rev{agent-swarm} world constrain
\rev{social-graph} information, spending, and settlement while making the
effects of alternative information and interaction rules inspectable?}
Applications include capital allocation, participation incentives, disclosure
policies, and credit stress scenarios in decentralized finance. Reports connect
financial outcomes \rev{and simulated price paths} with coordination,
\rev{social-graph} message delivery, and \rev{swarm} reactions to events so
researchers, protocol designers, and risk analysts can trace how an outcome
arises under declared assumptions.

Accordingly, we evaluate \textit{Boyun} through six synthetic mechanism
experiments, a fourteen-day \rev{sixteen-agent swarm} token-launch case across
24 policy paths and a ninety-day \rev{twenty-agent swarm} collateralized lending
case across four worlds---all executed locally with \textit{Laya}---together
with a matched-brief comparison against the original \textit{MiroFish} \rev{that
inspired \textit{Boyun}'s workflow}. These tests do not establish human
behavioral fidelity, forecasting accuracy, or improved user decisions. Within
that scope, the paper makes five contributions:
\begin{enumerate}\setlength{\itemsep}{1pt}
\item an explicit specification of worlds, \rev{the \textit{swarm of agents},
and what each agent may observe over the \textit{social graph}}
(Section~\ref{sec:world-creation});
\item an execution procedure across strategic, market, and lending ledgers that
records every \rev{agent} decision and replays it exactly
(Section~\ref{sec:method-system});
\item reproducible mechanism tests with known answers
(Sections~\ref{sec:evaluation}--\ref{sec:results});
\item a matched case comparison with the original \textit{MiroFish}
(Section~\ref{sec:res-mirofish});
\item full-horizon local \textit{Laya} token-launch and lending case studies
\rev{---including simulated token-price trajectories across 24 worlds---}that
separate mechanical conclusions fixed by ledger rules from behavioral ones
sensitive to the \rev{agent} decision model (Sections~\ref{sec:res-macron}
and~\ref{sec:res-lending}).
\end{enumerate}
The following section situates these contributions within prior work.
